\notechapter{N-20260311}{Matrix Computation I: Norms, Amplification, and Numerical Size}

\section{A matrix can be small entrywise and large as an operator}

A matrix is not merely a rectangular table of numbers.  It is a map
\[
  x\longmapsto Ax,
\]
and numerical analysis cares about what that map can do to perturbations.

Suppose the input is changed from \(x\) to \(x+\delta x\).  The output changes by
\[
  A(x+\delta x)-Ax=A\delta x.
\]
To say anything useful, we need a notion of size for both \(\delta x\) and \(A\delta x\).  A vector norm measures the size of data; an operator norm measures the largest factor by which a matrix can enlarge that size.

This distinction matters.  The matrix
\[
  A=
  \begin{bmatrix}
    1&1&\cdots&1\\
    1&1&\cdots&1\\
    \vdots&\vdots&\ddots&\vdots\\
    1&1&\cdots&1
  \end{bmatrix}
\]
has entries no larger than \(1\), yet in the all-ones direction it multiplies Euclidean length by \(n\).  Looking at the largest entry alone misses the collective action of the columns.

\section{Different norms ask different questions}

For \(x=(x_1,\ldots,x_n)^T\), the standard norms are
\[
  \|x\|_1=\sum_{i=1}^n|x_i|,
  \qquad
  \|x\|_2=\left(\sum_{i=1}^n|x_i|^2\right)^{1/2},
\]
\[
  \|x\|_\infty=\max_i|x_i|.
\]
More generally,
\[
  \|x\|_p=\left(\sum_{i=1}^n|x_i|^p\right)^{1/p},
  \qquad 1\le p<\infty.
\]

These norms emphasize different features.  The \(1\)-norm measures total mass.  The \(2\)-norm is rotationally invariant and fits inner-product geometry.  The \(\infty\)-norm records the worst coordinate.  For
\[
  x=(1,1,\ldots,1)^T,
\]
we have
\[
  \|x\|_1=n,
  \qquad
  \|x\|_2=\sqrt n,
  \qquad
  \|x\|_\infty=1.
\]
None of these answers is more correct in isolation.  They answer different questions about what counts as a large perturbation.

In finite dimensions all norms are equivalent: for any two norms \(\|\cdot\|_a\) and \(\|\cdot\|_b\), there exist constants \(c,C>0\) such that
\[
  c\|x\|_a\le \|x\|_b\le C\|x\|_a.
\]
Thus they define the same convergent sequences, but the constants may depend strongly on dimension.  Numerical estimates therefore do depend on the chosen norm even when qualitative convergence does not.

\section{Hölder duality explains the familiar inequalities}

Let \(1\le p,q\le\infty\) satisfy
\[
  \frac1p+\frac1q=1.
\]
Hölder's inequality states
\[
  |y^Tx|\le \|y\|_q\|x\|_p.
\]
The exponent \(q\) is called the conjugate exponent of \(p\).

This inequality says that a linear functional \(x\mapsto y^Tx\) has size \(\|y\|_q\) when inputs are measured in \(p\)-norm.  More precisely,
\[
  \max_{\|x\|_p=1}|y^Tx|=\|y\|_q.
\]
For \(1<p<\infty\), equality is achieved by choosing the phases of \(x_i\) to match \(y_i\) and setting
\[
  |x_i|\propto |y_i|^{q-1}.
\]
For \(p=1\), the maximizing vector concentrates on a coordinate where \(|y_i|\) is largest; hence the dual of \(\ell^1\) is \(\ell^\infty\).  For \(p=\infty\), choose each coordinate with maximal allowed magnitude and matching sign; the dual is \(\ell^1\).

The Cauchy--Schwarz inequality is simply the case \(p=q=2\):
\[
  |y^Tx|\le \|y\|_2\|x\|_2.
\]
This dual viewpoint is useful because matrix norms can be read as collections of linear functionals acting on columns or rows.

\section{An induced norm measures worst-case amplification}

Given norms on the input and output spaces, define
\[
  \|A\|=\max_{x\ne0}\frac{\|Ax\|}{\|x\|}
  =\max_{\|x\|=1}\|Ax\|.
\]
This is the induced, or subordinate, operator norm.  By definition,
\[
  \|Ax\|\le \|A\|\,\|x\|.
\]
It is the smallest constant for which this inequality holds for every \(x\).

Induced norms are automatically submultiplicative:
\[
  \|AB\|
  =\max_{x\ne0}\frac{\|ABx\|}{\|x\|}
  \le
  \max_{x\ne0}
  \frac{\|A\|\,\|Bx\|}{\|x\|}
  \le \|A\|\|B\|.
\]
Thus a product of linear maps cannot amplify more than the product of their individual worst-case factors.

The standard formulas are
\[
  \|A\|_1=\max_j\sum_i|a_{ij}|,
  \qquad
  \|A\|_\infty=\max_i\sum_j|a_{ij}|.
\]
The first is the largest absolute column sum; the second is the largest absolute row sum.

To see the \(1\)-norm formula, write
\[
  \|Ax\|_1
  \le \sum_i\sum_j|a_{ij}|\,|x_j|
  =\sum_j\left(\sum_i|a_{ij}|\right)|x_j|.
\]
This is bounded by
\[
  \left(\max_j\sum_i|a_{ij}|\right)\|x\|_1.
\]
Equality is obtained by putting all mass on a column with maximal sum and choosing the sign or phase appropriately.

\section{A worked comparison of three operator norms}

Let
\[
  A=
  \begin{bmatrix}
    1&-2\\
    2&2
  \end{bmatrix}.
\]
The absolute column sums are
\[
  3,\qquad4,
\]
so
\[
  \|A\|_1=4.
\]
The absolute row sums are
\[
  3,\qquad4,
\]
so
\[
  \|A\|_\infty=4.
\]

For the Euclidean norm, compute
\[
  A^TA=
  \begin{bmatrix}
    5&2\\
    2&8
  \end{bmatrix}.
\]
Its eigenvalues are
\[
  \lambda_\pm=\frac{13\pm5}{2}=9,4.
\]
Therefore
\[
  \|A\|_2=\sqrt{\lambda_{\max}(A^TA)}=3.
\]
The three values differ because they optimize over different unit balls.  The \(1\)- and \(\infty\)-balls have corners aligned with coordinate directions; the Euclidean ball is rotationally symmetric.

The right singular vector corresponding to singular value \(3\) is a direction in which \(A\) stretches Euclidean length exactly by a factor of \(3\).  Other input directions stretch less.

\section{The spectral norm comes from singular values}

The Euclidean operator norm satisfies
\[
  \|A\|_2
  =\max_{\|x\|_2=1}\|Ax\|_2.
\]
Squaring gives
\[
  \|Ax\|_2^2=x^TA^TAx.
\]
The maximum of this Rayleigh quotient over the unit sphere is the largest eigenvalue of \(A^TA\).  Hence
\[
  \boxed{
  \|A\|_2=\sqrt{\lambda_{\max}(A^TA)}=\sigma_{\max}(A).}
\]

This formula explains why singular values, rather than eigenvalues, govern Euclidean amplification.  Eigenvalues describe invariant directions satisfying \(Ax=\lambda x\).  Singular vectors solve a different optimization problem: which input direction produces the longest output?

For a nonnormal matrix these questions can diverge dramatically.  Consider
\[
  N_M=
  \begin{bmatrix}
    0&M\\
    0&0
  \end{bmatrix}.
\]
Both eigenvalues are zero, so the spectral radius is zero, yet
\[
  \|N_M\|_2=M.
\]
The matrix is nilpotent, but it can cause a large one-step transient amplification.  This is why eigenvalues alone are not a complete measure of numerical size.

\section{Unitary transformations preserve Euclidean geometry}

If \(U^TU=I\), then
\[
  \|Ux\|_2^2=x^TU^TUx=\|x\|_2^2.
\]
Thus orthogonal or unitary matrices are changes of orthonormal coordinates.  They rotate or reflect without changing length.

Consequently,
\[
  \|UAV\|_2=\|A\|_2
\]
for unitary \(U,V\).  Singular value decomposition writes
\[
  A=U\Sigma V^T,
\]
so the Euclidean operator norm is simply the largest diagonal entry of \(\Sigma\).

The Frobenius norm
\[
  \|A\|_F=\left(\sum_{i,j}|a_{ij}|^2\right)^{1/2}
\]
is also unitarily invariant.  Since
\[
  \|A\|_F^2=\operatorname{tr}(A^TA),
\]
we obtain
\[
  \|UAV\|_F^2
  =\operatorname{tr}(V^TA^TU^TUAV)
  =\operatorname{tr}(V^TA^TAV)
  =\operatorname{tr}(A^TA).
\]
Equivalently,
\[
  \|A\|_F^2=\sum_i\sigma_i(A)^2.
\]

The spectral norm asks for the single strongest direction.  The Frobenius norm combines energy across all singular directions.  They satisfy
\[
  \|A\|_2\le \|A\|_F\le \sqrt{r}\,\|A\|_2,
\]
where \(r=\operatorname{rank}A\).

\section{Schatten norms place spectral and Frobenius size in one family}

For \(1\le p<\infty\), define the Schatten \(p\)-norm by applying a vector \(p\)-norm to the singular values:
\[
  \|A\|_{S_p}
  =\left(\sum_i\sigma_i(A)^p\right)^{1/p}.
\]
Then
\[
  \|A\|_{S_2}=\|A\|_F,
  \qquad
  \|A\|_{S_\infty}=\|A\|_2.
\]
The case \(p=1\) is the nuclear norm,
\[
  \|A\|_*=\sum_i\sigma_i(A).
\]
It is used as a convex surrogate for rank because it penalizes the total singular-value mass.

These norms separate two notions that are often conflated.  Entrywise sparsity concerns how many coordinates are nonzero.  Low rank concerns how many singular directions are active.  The \(1\)-norm of entries encourages the former; the nuclear norm encourages the latter.

\section{Norms turn perturbation statements into inequalities}

If \(x\) is perturbed by \(\delta x\), then
\[
  \|A(x+\delta x)-Ax\|
  \le \|A\|\,\|\delta x\|.
\]
This is the basic Lipschitz estimate for a linear map.

For a product,
\[
  F(x)=W_LW_{L-1}\cdots W_1x,
\]
we have
\[
  \|F(x+\delta x)-F(x)\|_2
  \le
  \left(\prod_{\ell=1}^L\|W_\ell\|_2\right)
  \|\delta x\|_2.
\]
The bound can be pessimistic because the maximally amplified direction of one layer may not align with that of the next.  Still, it gives a rigorous worst-case control and explains why spectral normalization constrains the Lipschitz constant of a network layer.

Norms also bound truncation and rounding errors.  If a computed matrix is \(A+E\), then
\[
  \|(A+E)x-Ax\|
  \le \|E\|\,\|x\|.
\]
The norm of the perturbation is meaningful only relative to the chosen geometry.  A small Frobenius perturbation controls average singular energy, while a small spectral perturbation controls every Euclidean input direction.

\section{The distance to singularity appears in the smallest singular value}

For an invertible square matrix,
\[
  \sigma_{\min}(A)
  =\min_{\|x\|_2=1}\|Ax\|_2.
\]
This is the least amount by which \(A\) stretches any unit direction.  If it is small, there is a direction that is almost collapsed.

The Euclidean distance from \(A\) to the set of singular matrices is exactly
\[
  \sigma_{\min}(A)
\]
when distance is measured by the spectral norm.  Indeed, if \(u_n,v_n\) are singular vectors for \(\sigma_{\min}\), then
\[
  E=-\sigma_{\min}u_nv_n^T
\]
has norm \(\sigma_{\min}\) and makes
\[
  (A+E)v_n=0.
\]
No smaller perturbation can make the matrix singular, because
\[
  \|(A+E)x\|_2
  \ge \|Ax\|_2-\|E\|_2\|x\|_2
  \ge \sigma_{\min}(A)-\|E\|_2
\]
for unit \(x\).

Thus the smallest singular value is not merely one more number in an SVD table.  It measures how close invertibility is to failure.

\section{A norm is useful only when its role is stated}

Several quantities called ``matrix size'' coexist because they serve different purposes.

The maximum entry answers a local coordinate question.  The Frobenius norm records total entrywise or singular-value energy.  An induced norm controls worst-case amplification for a chosen notion of vector size.  The spectral norm is the Euclidean Lipschitz constant.  The nuclear norm measures total singular mass and is sensitive to low-rank structure.  The smallest singular value measures the weakest direction and the distance to singularity.

The practical habit should therefore be to complete the sentence:
\[
  \text{``This norm is being used to control \emph{what}?''}
\]
Without that clause, a numerical bound may be formally correct but conceptually irrelevant.  With it, norms become the language that connects geometry, perturbations, algorithms, and stability.
